\subsection{THE COMPLEX EXPONENTIAL}

We have determined (Gen.~Top., VIII, p.~106) all the continuous homomorphisms of the (additive) topological group C of complex numbers onto the (multiplicative) topological group \(C^{*}\) of complex numbers \(\neq 0\) ; these are the maps

\[
x + i y \mapsto e ^ {\alpha x + \beta y} \mathbf {e} (\gamma x + \delta y)\tag{20}
\]

where \(\alpha, \beta, \gamma, \delta\) are four real numbers subject to the single condition \(\alpha\delta - \beta\gamma \neq 0\) . We now propose to determine which of these homomorphisms \(z \mapsto f(z)\) are differentiable on C. First we remark that it is enough for f to be differentiable at the point z = 0; indeed, for every point \(z \in C\) one has \(\frac{f(z + h) - f(z)}{h} = f(z) \frac{f(h) - 1}{h}\) ; if \(f'(0)\) exists, then so does \(f'(z)\) , and \(f'(z) = af(z)\) , with \(a = f'(0)\) . On the other hand, if g is a second differentiable homomorphism, such that \(g'(z) = bg(z)\) , then \(g(az/b) = f(z)\) , for one notes immediately that the quotient \(g(az/b)/f(z)\) has an everywhere zero derivative and is equal to 1 for z = 0; all the differentiable homomorphisms are thus of the form \(z \mapsto f(\lambda z)\) , where f is one of them (assuming they exist) and \(\lambda\) is any (complex) constant.

This being so, if \(f\) is differentiable at the point \(z = 0\) then each of the maps \(x \mapsto f(x)\) , \(y \mapsto f(iy)\) of \(\mathbf{R}\) into \(\mathbf{C}\) is necessarily differentiable at the point 0, the first having derivative \(f'(0)\) , the second \(if'(0)\) . Now the derivatives of the maps \(x \mapsto e^{\alpha x}\mathbf{e}(\gamma x)\) , \(y \mapsto e^{\beta y}\mathbf{e}(\delta y)\) at the point 0 are respectively equal to \(\alpha + 2\pi i\gamma\) and \(\beta + 2\pi i\delta\) , from which \(\beta = -2\pi\gamma\) and \(\alpha = 2\pi\delta\) ; these conditions are, in particular, satisfied by the homomorphism \(x + iy \mapsto e^x\mathbf{e}(y/2\pi)\) , which we shall denote provisionally by \(f_0\) . We shall now show that \(f_0\) is actually differentiable at the point \(z = 0\) .

It is clear that \(x \mapsto f_0(x)\) and \(y \mapsto f_0(iy)\) have derivatives of every order; in particular, Taylor's formula of order 1 applied to these functions shows that for every \(\varepsilon > 0\) there is an \(r > 0\) such that, if one puts

\[
f _ {0} (x) = 1 + x + \varphi (x) x, \quad f _ {0} (i y) = 1 + i y + \psi (y) y,
\]

then the conditions \(|x| \leqslant r, |y| \leqslant r\) imply that \(|\varphi(x)| \leqslant \varepsilon\) and \(|\psi(y)| \leqslant \varepsilon\) ; this being so, we have \(f_0(x + iy) = f_0(x)f_0(iy) = 1 + (x + iy) + \theta(x, y)\) with

\[
\theta (x, y) = \bigl (i + \varphi (x) \psi (y) \bigr) x y + (1 + x) y \psi (y) + (1 + i y) x \varphi (x);
\]

for \(|z| \leqslant r\) we have \(|x| \leqslant r\) and \(|y| \leqslant r\) , whence

\[
\left| \theta (x, y) \right| \leqslant \left(1 + \varepsilon^ {2}\right) | z | ^ {2} + 2 \varepsilon | z | (1 + | z |)
\]

which proves that the quotient \(\frac{f_{0}(z)-1-z}{z}\) tends to 0 with z, that is to say, the function \(f_{0}\) admits a derivative equal to 1 at the point z=0. The above thus proves that, for all \(z\in C\) ,

\[
\mathrm{D} \big (f _ {0} (z) \big) = f _ {0} (z).\tag{21}
\]

This property establishes the connection between \(f_{0}\) and the function \(e^{x}\) , which is moreover the restriction of \(f_{0}\) to the real axis; for this reason we make the following definition:

DEFINITION 2. The homomorphism \(x + iy \mapsto e^x\mathbf{e}(y/2\pi)\) of \(\mathbf{C}\) onto \(\mathbf{C}^*\) is called the complex exponential; its value at an arbitrary complex number \(z\) is denoted by \(e^z\) or \(\exp z\) .

